\section{Discussion and limitations}

The main conceptual lesson is representational.  Spectral normalization,
local Gaussian structure, logarithmic differentiation, and cut-locus
multiplicity are easiest in different coordinates.  A rigorous calculation
therefore needs a domain atlas and explicit transfer conditions rather than a
single globally preferred formula.

The constants in Table~\ref{tab:certificate} are intentionally conservative.
Replacing the dyadic bound $e^{-x}\le2^{-\lfloor x\rfloor}$ with rational
range-reduced Taylor bounds improves them without changing the theorem.  A
branch-and-bound interval implementation could also subdivide $t$ and
$\theta$ to reduce dependency.  Sharpness is secondary here to a short
collection of elementary inequalities.

The global theorem separates two genuinely different regimes.  After removal
of the Gaussian action, the identity-chart image correction is exponentially
small uniformly through any requested finite order.  At the antipode the
$S^2$-family of minimizing geodesics changes the heat-kernel prefactor and
forces stronger growth of logarithmic derivatives.  The cut-chart constants
therefore retain their $t$ dependence.  This is a geometric feature, not a
loss introduced by the proof.

Several extensions are natural but not proved.  The Weyl denominator and
affine Weyl group suggest analogues for higher-rank compact groups, but walls
intersect and the endpoint atlas becomes multidimensional.  A result for
$\mathrm{SU}(N)$ would require controlling those intersections, not simply
replacing $2\pi k$ by a vector.  Likewise, the local action budgets become an
RG input only after dependent-measure, large-field, gauge, and topology
estimates are supplied.
